% =====================================================================
% Block M1, v3 (theory/v3): shared assumptions, the deployed default (def:m1-default), decoder target, where the
% input enters, the autonomous case (prop:m1-autonomous), the non-autonomous residual (rem:m1-tv).
% v3 changes are marked \vthree{ID}{...}; IDs in theory/v3/STATUS.md.
% Owns: Assumption ass:m1 (A1)-(A4) used by every block, eq:m1-TL, def:m1-suff, lem:m1-window,
% prop:m1-suff, rem:m1-koopman, rem:m1-capacity, prop:m1-where, cor:m1-amplitude, prop:m1-tight, cor:m1-ood.
% Cites: lem:m2-tail, thm:m2-window, prop:m2-current, lem:m2-sfm (block M2); lem:D, prop:m3-iss, prop:m3-mismatch (block M3).
% Checks: paragraph "Numerical checks" at the end of this section (theory/checks/m1_location.py, check_theory.py).
% Status legend used in all three sections: [proven] / [proven, linear plants] / [conjecture] / [empirical].
% =====================================================================
\section{Where the input enters a learned KKL observer}\label{sec:m1}

\paragraph{Setting and conventions.} Plant $\dot x=f(x,u)$, $y=h(x)$, $x\in\mathbb R^{n_x}$,
$y\in\mathbb R^{n_y}$; no input-affine structure. Inputs are applied through a zero-order hold at the
sample rate with values in a compact $U\ni0$, $\bar u:=\max_{U}|u|$; $u_t$ is the window $u|_{[t-L,t]}$
shifted to $[-L,0]$, a point of $W_L=U^{N_L}$ at sample instants when $L/dt=N_L\in\mathbb N$ (the decoder
runs only there), with the sup norm. $X(x,t;s;u)$ is the value at $s$ of the solution through $x$ at $t$.
The latent filter is $\dot z=Az+By$. \emph{Latent initial condition (one convention for Sections
\ref{sec:m1}--\ref{sec:m3}).} The \emph{label} and training latents start from $z(0)=0$, so
$z(t)=T^0(x(t),t;u):=\int_0^te^{A(t-s)}Bh(X(x(t),t;s;u))\,ds$ exactly; this finite-start map is the only
time-varying map used. The \emph{deployed} observer (decoder-side and latent-injection alike, \texttt{evaluate.py}
\texttt{z0\_map}) starts from $\hat z(0)=\hat T_0(\hat x_0)$, which adds the transient $e^{At}\hat T_0(\hat x_0)$ to
the latent; it is carried as the $\Delta z_0$ term of Proposition~\ref{prop:m3-iss} and never absorbed into a map.
The decoder target is the windowed map
\begin{equation}\label{eq:m1-TL}
T_L(x,w):=\int_{-L}^{0}e^{-A\tau}B\,h\big(X(x,0;\tau;w)\big)\,d\tau,
\end{equation}
which by Lemma~\ref{lem:m2-tail}(a) is within $\tfrac12\varepsilon(L)$ of $T^0(x(t),t;u)$ at $w=u_t$.
$\mathcal P_L$ is the set of pairs $(x(t),u_t)$, $t\ge L$, generated by admissible trajectories.

\begin{definition}[the deployed default observer]\label{def:m1-default}
\vthree{M1-1}{new: the final method, A3 + mixed histories + C5} Latent filter (fixed, never trained):
$\hat z_{k+1}=A_d\hat z_k+W_0\tilde y_k+W_1\tilde y_{k+1}$, $\hat z_0=\hat T_0(\hat x_0)$ (Lemma~\ref{lem:D}(a)).
Latent window: $Z_k:=(\hat z_k,\hat z_{k-s},\dots,\hat z_{k-(K-1)s})$, entries with negative index set to $0$, $K=20$,
stride $s=\max\{1,\mathrm{round}(6/(\alpha\,dt\,K))\}$, span $S:=(K-1)s\,dt$; on the twelve benchmark plants
$\alpha S\in[5.50,5.90]$ (check V8). Input context (SFM): $\xi_{k+1}=A^\xi_d\xi_k+B_du_k/u_{\rm ref}$, $\xi_0=0$,
$c_k=\psi([\xi_k,u_k/u_{\rm ref}])$ with $\psi$ a bias-free $\tanh$ MLP. Decoder:
\[
\hat x_k=D(Z_k,c_k):=T^+\big(\hat z_k;\,[c_k,g(Z_k)]\big),
\]
$g$ a bias-free $\tanh$ MLP of the normalised stack, $T^+$ the base MLP (activation $\tanh$ in the stored default
models, check V7) whose layer $l$ is modulated by $\Delta W_l(t)=U_l(t)V_l(t)/\sqrt r$, $t=\mathrm{Trunk}([c,g])$
(bias-free SiLU), $U_l$ affine, $V_l$ linear, $r=4$, with the spectral cap $M\mapsto M\min\{1,\gamma/\|M\|_2\}$ on
every head weight. The base decoder is retrained in phase~2. Training uses the mixed histories $D_0\cup D_1$ and the
loss with the C5 terms (jump penalty and off-manifold slope penalties); these change what the decoder is fitted to, not
the structure above. \emph{Not in the default:} an input-dependent latent gain $B(c)$ and a latent injection $\Phi$.
\emph{Special case $K=1$} (no $z$-window, frozen base): $D(Z,c)=T^+(\hat z_k;c)$; this is the observer of the C1 and C2
comparison nodes (recipe A0), and every statement below that is written for $D(z,c)$ refers to it.
\end{definition}
\noindent Two exact consequences used below. (i) The stack enters only through $D$, so every result that bounds
the latent deviation (Proposition~\ref{prop:m3-iss}(i)) is unchanged; what changes is the Jacobian with respect to
$Z$ (Proposition~\ref{prop:m3-iss}). (ii) The spectral cap gives, for every context,
$\|\Delta W_l(t)\|_2\le\|U_l(t)\|_F\|V_l(t)\|_F/\sqrt r\le(\gamma\|t\|+\|b_{U_l}\|)\gamma\|t\|/\sqrt r$ ($b_{U_l}$ the
head bias), checked on the stored default models (V7c: ratio $\le0.67$). It bounds the head part of the context
Lipschitz constant on bounded sets; the trunk and base layers still enter as norm products, so it does not make
Proposition~\ref{prop:m3-iss}(iii) informative a priori. \vthree{M1-2}{new}

\begin{assumption}[standing; shared by all sections]\label{ass:m1}
Regularity: $f(\cdot,u)\in C^1$ uniformly in $u\in U$, $h\in C^2$.
\begin{enumerate}
\item[(A1)] \emph{Forward invariance.} A compact $\mathcal X$ contains every solution from the initial set
  $\mathcal X_0\subset\mathcal X$ under every admissible input, for all $t\ge0$.
\item[(A2)] \emph{Bounded output.} $\|h(x)\|\le\bar y$ on $\mathcal X$.
\item[(A3)] \emph{Hurwitz filter.} $-\alpha<0$ is the spectral abscissa of $A$ and
  $\kappa:=\sup_{t\ge0}e^{\alpha t}\|e^{At}\|<\infty$ ($\kappa=1$ for the normal Butterworth $A$ of the code).
\item[(A4)] \emph{Uniform Lipschitz injectivity of the windowed map.} There are $L_0$ and $\ell<\infty$ with
  $|x_1-x_2|\le\ell\,|T_L(x_1,w)-T_L(x_2,w)|$ for all $(x_1,w),(x_2,w)\in\overline{\mathcal P_L}$ and all
  $L\ge L_0$ (i.e.\ $\ell:=\sup_{L\ge L_0}\ell(L)<\infty$; Theorem~\ref{thm:m2-window}(a) varies $L$ at fixed $\ell$).
\end{enumerate}
\end{assumption}
\noindent\emph{Chaotic plants.} (A1)--(A3) hold for the dissipative chaotic benchmarks (Lorenz, R\"ossler), (A4) is
assumed for them as for every other plant, but its constant is not under control when $\alpha$ is below the
\emph{backward} expansion rate $\lambda_b$ of the flow (for Lorenz $\lambda_b\approx14.6$ against $\alpha=2$):
$\partial_xT^0$ contains $\int_0^te^{A(t-s)}B\,\partial h\,\partial_xX(x,t;s;u)\,ds$ and $\|\partial_xX\|$ grows like
$e^{\lambda_b(t-s)}$ along the backward flow, so the forward Lipschitz constant of $T^0(\cdot,t)$ grows with $t$ when
$\alpha<\lambda_b$; this is the eigenvalue condition of \citet{andrieu2006existence} (the spectrum of $A$ must
dominate the Lipschitz constant of $f$ for a $C^1$ stationary map; a H\"older map needs less), and it is a plausible
mechanism (not tested: no node sweeps $\alpha$ or measures $\lambda_b$) behind the phase-1 PDE residual
$\varepsilon_{\rm pde,rms}=74$ on Lorenz against $\le0.3$ elsewhere. Which results survive on such plants is stated at
each result: everything that uses only the true past and (A1)--(A2) does (Lemma~\ref{lem:m2-tail}(a),
Theorem~\ref{thm:m2-window}(a), Proposition~\ref{prop:m3-iss}); the counterfactual-past statements
(Lemma~\ref{lem:m2-tail}(b)) and the linear lower bounds (Theorem~\ref{thm:m2-window}(b),
Proposition~\ref{prop:m2-current}(b)) do not apply to them.

\begin{remark}[existence and injectivity: cited, and what is assumed]\label{rem:m1-exist}
$T^0$ is a $C^1$ solution of $\partial_xT\,f(x,u(t))+\partial_tT=AT+Bh(x)$ whenever backward solutions
exist on $[0,t]$ \citep[Lemma~1]{bernard2019luenberger}, and an observer
$\hat x=T^*(\hat z,t)$ converges if there are $\bar t$ and a concave class-$\mathcal K_\infty$ $\rho_T$ with
$|x_1-x_2|\le\rho_T(|T(x_1,t)-T(x_2,t)|)$ for all $t\ge\bar t$ \citep[Theorem~1]{bernard2019luenberger}.
For a fixed input that is backward distinguishable in finite time, $T^0(\cdot,t)$ is injective for almost
every complex diagonal $A$ with $n_y(n_x+1)$ eigenvalues \citep[Theorem~3]{bernard2019luenberger}, but the
exceptional set depends on the input \citep[Remark~5]{bernard2019luenberger}; the Lipschitz, uniform-in-time
version \citep[Theorem~2]{bernard2019luenberger} needs strong differential observability, inputs of class
$C^{m+1}$ (their time derivatives enter) and a gain $k$ above a threshold, which excludes held, switching
and PRBS inputs and a fixed $(A,B)$. For autonomous systems, $2n_x+1$ generic real dimensions suffice
\citep[Theorem~3.4]{brivadis2023further}, $2n_x+2$ for time-varying ones \citep[Remark~3.6]{brivadis2023further}.
(A4) is therefore \emph{assumed}: it is the quantitative, uniform-in-input version of what these results give
for each regular input. The code's $n_z=5$ for 2-D plants is the autonomous count; the time-varying count
$2n_x+2$ is not met, so (A4) at $n_z=5$ is an assumption whose constant must be reported, not a consequence of the cited
results; the phase-1 proxy for $\ell$ (\texttt{l\_dec\_q95}, Remark~\ref{rem:m3-cert}) is $24$ (vdp), $44$
(bioreactor), $299$ (duffing), $10$ (lorenz), $14$ (quad\_tank), $816$ (battery) for the C1-where phase~1 (used in
\texttt{check\_apriori\_vacuity}); the C2-window-a1 phase~1 of the results layer gives $23,44,295,10,14,806$
(Table~\ref{tab:m2-window-rule}).
\end{remark}

\begin{remark}[Koopman reading; which model classes the variants are]\label{rem:m1-koopman}
For the autonomous plant and a diagonalised filter $A=\mathrm{diag}(\lambda_i)$, $B=(b_i)$, the KKL PDE
$\partial_xT_i\,f=\lambda_iT_i+b_ih$ reads $(L_f-\lambda_i)T_i=b_ih$ with $L_f=f\cdot\nabla$ the Koopman generator on
observables, so
\[
T_i=b_i\,(L_f-\lambda_i)^{-1}h=b_i\int_0^\infty e^{\lambda_i\tau}\,h\big(X(x,0;-\tau;0)\big)\,d\tau ,
\]
the resolvent of the generator at the filter pole, i.e.\ the Laplace transform of the backward Koopman semigroup
applied to the output observable; if $h=\sum_jc_j\varphi_j$ with $L_f\varphi_j=\mu_j\varphi_j$ then
$T_i=\sum_jb_ic_j\varphi_j/(\mu_j-\lambda_i)$. For a linear plant this is the row identity
$p_i=b_iC(F-\lambda_iI)^{-1}$ of the Sylvester solution (check [7]); for a nonlinear plant it gives the spectral
reading of the ``$A$ fast enough'' condition above ($\operatorname{Re}\lambda_i$ must clear the backward spectrum) and of
the zero-input floor $\varepsilon_0$ of Remark~\ref{rem:m1-capacity}(i) as a spectral truncation. Under an input the
same object is the resolvent for the skew-product (state $\times$ input-shift) system, and the windowed map $T_L$ is its
finite-horizon version with the window as a truncated coordinate on the shift space; for a linear plant the context
$\dot q=Aq-PGu$ of Proposition~\ref{prop:m1-suff}(a) has the structure of a linear predictor with inputs
\citep{proctor2016dmdc,korda2018linear}. In these terms, \emph{latent injection} with a stationary $T_0$ is a lifted
Koopman model with inputs (its input-affine case $\Omega(z)u$, and the NLOX baseline $\omega(z)u$, are bilinear lifted
observers with a fixed output-injection gain $B$), whose stability requires the input-dependent generator $A+\partial_z\Phi$
to stay contracting (Proposition~\ref{prop:m1-where}(a)); \emph{decoder-side conditioning} keeps a fixed generator and
parametrises the left inverse of the lifting by the input context. Proposition~\ref{prop:m1-tight} is then the
statement that the Koopman spectrum moves with $u$ (eigenvalue $u-a$; rotation rate $\omega_0+u$) while a fixed $A$ does
not. We state this as positioning only; no result below depends on it.
\end{remark}

\subsection{The decoder target and sufficient contexts}

\begin{definition}[sufficient context]\label{def:m1-suff}
A context map $E:W_L\to\mathbb R^{d_c}$ is \emph{sufficient} if the latent and the context determine the state:
for $(x_i,w_i)\in\overline{\mathcal P_L}$, $E(w_1)=E(w_2)\Rightarrow|x_1-x_2|\le\ell\,|T_L(x_1,w_1)-T_L(x_2,w_2)|$.
It is $\epsilon_E$-sufficient if the same holds with an additive $\epsilon_E$ on the right.
\end{definition}
\noindent The window itself ($E=\mathrm{id}$) is sufficient by (A4). A lower-dimensional $E$ cannot be
injective on $W_L$ (invariance of domain), but it need not be: it only has to resolve what the latent leaves open.

\begin{lemma}[decoder target]\label{lem:m1-window}
Let $E$ be continuous and $\epsilon_E$-sufficient, and $S_L:=\{(T_L(x,w),E(w)):(x,w)\in\overline{\mathcal P_L}\}$.
If $\epsilon_E=0$ there is a continuous $G^*:S_L\to\mathcal X$ with $G^*(T_L(x,w),E(w))=x$, $\ell$-Lipschitz in $z$
for each $c$; extend $G^*(\cdot,c)$ to $\mathbb R^{n_z}$ with the same constant (Kirszbraun). Then along every
admissible trajectory, for $t\ge L$,
$\ |G^*(z(t),E(u_t))-x(t)|\le\tfrac{\ell}{2}\varepsilon(L).$
For $\epsilon_E>0$, any decoder $D$ satisfies $\sup|D(T_L(x,w),E(w))-x|\ge\tfrac12\delta_E$, where $\delta_E$ is the
largest $|x_1-x_2|$ over pairs with equal latent and equal context.
\end{lemma}
\noindent\emph{Intuition.} Over a fixed window the plant has no clock, so the state is a fixed function of
(latent, context) as soon as the context resolves the latent's ambiguity.
\begin{proof}
$\Psi(x,w)=(T_L(x,w),E(w))$ is continuous on the compact $\overline{\mathcal P_L}$ (under the hold, the backward
solution is a composition of $N_L$ $C^1$ flows with constant parameters). By sufficiency, $\Psi(x_1,w_1)=\Psi(x_2,w_2)$
implies $x_1=x_2$, so $\pi_x$ factors as $G^*\circ\Psi$; $G^*$ is continuous because $\Psi$ is a closed map
(continuous on a compact set into a Hausdorff space), and $\ell$-Lipschitz in $z$ by sufficiency. Then
$|G^*(z(t),c)-x(t)|=|G^*(z(t),c)-G^*(T_L(x(t),u_t),c)|\le\ell\,|z(t)-T_L(x(t),u_t)|$ and Lemma~\ref{lem:m2-tail}(a).
The lower bound: $D$ returns one value for two states $\delta_E$ apart.
\end{proof}

\begin{proposition}[sufficient contexts exist, and must carry history]\label{prop:m1-suff}
(a) \emph{Linear plants.} For $\dot x=Fx+Gu$, $y=Cx$, the windowed map is affine in $x$ with a
$w$-independent slope, $T_L(x,w)=P_Lx+Q_Lw$, $P_L=P-e^{AL}Pe^{-FL}$, $PF=AP+BC$. Hence the $n_z$-dimensional linear
context $c=Q_Lw$ is exactly sufficient whenever $P_L$ is injective ($\ell=\|P_L^+\|$), and as $L\to\infty$ (under the
washout $e^{AL}Pe^{-FL}\to0$, (H1) of Theorem~\ref{thm:m2-window}) it is the
output $q$ of the stable filter $\dot q=Aq-PGu$: a spectral filter bank with identity readout.
(b) \emph{Nonlinear plants.} The window ($d_c=n_uN_L$) is sufficient, and by (a) of Lemma~\ref{lem:m2-tail}
any context that determines $u_t$ for some $L$ is $\ell\varepsilon(L)$-sufficient for the infinite past:
$|x_1-x_2|\le\ell\,|T^0_1-T^0_2|+\ell\,\varepsilon(L)$ (each $T^0_i$ is within $\tfrac12\varepsilon(L)$ of $T_L$).
(c) \emph{History is necessary.} The current input alone is not sufficient in general
(Proposition~\ref{prop:m2-current}).
\end{proposition}
\begin{proof}
(a) $X(x,0;\tau;w)=e^{F\tau}x+(\text{linear in }w)$; substitute in \eqref{eq:m1-TL} and use
$\int_{-L}^0e^{-A\tau}BCe^{F\tau}d\tau=P-e^{AL}Pe^{-FL}$ (integrate $\tfrac{d}{d\tau}(e^{-A\tau}Pe^{F\tau})=e^{-A\tau}BCe^{F\tau}$).
If $Q_Lw_1=Q_Lw_2$ then $T_L(x_1,w_1)-T_L(x_2,w_2)=P_L(x_1-x_2)$. The limit: $z-Px$ obeys $\dot{(z-Px)}=A(z-Px)-PGu$.
(b) (A4) and Lemma~\ref{lem:m2-tail}(a).
\end{proof}
\noindent For nonlinear plants and contexts of fixed dimension below $n_uN_L$, finite stable filter banks with a
static readout uniformly approximate every time-invariant fading-memory operator on slew-limited bounded inputs
\citep[Theorem~2]{boyd1985fading}; applying it to the observer operator $(y,u)\mapsto x$ (a fading-memory operator
under (A1)--(A4), on slew-limited signals, and for the filters the theorem chooses rather than ours) motivates the
spectral context but is not a guarantee for the trained encoder.

\begin{remark}[expressivity at a fixed base; what C1 compares]\label{rem:m1-capacity}
(i) \emph{Shared floor} \vthree{M1-3}{scope: the C1 arms, not the default}. Every conditioner with exact zero-input collapse, $D(\cdot,0)=T_0^+$, errs at $c=0$ by
$\varepsilon_0:=\sup_z|G^*(z,0)-T_0^+(z)|$: the zero-input target is the inverse of the zero-input windowed map,
not the learned stationary base. This describes the C1 comparison arms ($K=1$, frozen base). It does not describe the
default: there $D(\cdot,[0,g(Z)])\ne T_0^+$ and the base is retrained (Proposition~\ref{prop:m1-autonomous}(a)). (ii) \emph{Additive (concat).} $T_0^+(z)+R(z,c)-R(z,0)$ approximates on $S_L$, as
the width of $R$ grows, every continuous target that equals $T_0^+$ on the zero slice: SiLU satisfies
$\mathrm{SiLU}(s)-\mathrm{SiLU}(-s)=s$, so a two-hidden-layer SiLU network contains every one-hidden-layer one, and
the universal approximation theorem for non-polynomial activations \citep{leshno1993multilayer} applies (with a
Tietze extension of the target). (iii) \emph{Base-capped (FiLM, LoRA, hypernetwork).} For each fixed $c$ the decoder
is a network $N_{\theta(c)}$ of the frozen base's architecture, so for every conditioner budget
\[
\inf_{D}\ \sup_{S_L}|D-G^*|\ \ge\ \sup_{c}\ \inf_{\theta\in\Theta}\ \sup_z|N_\theta(z)-G^*(z,c)|,
\]
with $\Theta$ the reachable parameter set. At the matched budget $\Theta$ is much smaller than the base's weight
space: the hypernetwork's linear bias-free heads place $\theta(c)$ in an affine subspace of dimension equal to its
trunk width; LoRA updates have the configured rank and leave biases fixed; FiLM never modulates the readout layer, so
$D(\cdot,c)=W_L\phi_c(z)+b_L$ with $(W_L,b_L)$ frozen. In the ReLU setting, a LoRA rank of order width times a depth
ratio suffices to represent a smaller target network exactly \citep{zeng2024expressive}; that result does not cover
tanh bases or small fixed ranks. (iv) \emph{Latent injection.} $\Phi$ is a universal class, but as an observer it
represents only targets of the stationary route: a stationary $T_0$ with an exact correction.
\vthree{M1-4}{default} For the default the base weights are retrained in phase~2, so $\Theta$ also contains the
base's own weights; the display stays valid with that larger $\Theta$. The theory therefore does not rank the variants by expressivity; at a matched budget on one base, C1 measures
inductive bias and optimisation within these classes.
\end{remark}

\subsection{Where the input enters: a stability--extrapolation trade-off}

Both observers run the filter on $y+\nu$, $\|\nu\|_\infty\le\delta_y$, with a context $\hat c=E(\hat u)$ of the input
$\hat u$ they are given; $c=E(u)$ and $\delta_c:=\sup_t|\hat c-c|\le\ell_E\|\hat u-u\|_\infty$ (Lemma~\ref{lem:m2-sfm}).
\begin{align*}
\text{(latent)}\quad &\dot{\hat z}=A\hat z+B(y+\nu)+\Phi(\hat z,\hat c), && \hat x=T_0^+(\hat z),\ \ \hat z(0)=T_0(\hat x_0),\\
\text{(decoder side)}\quad &\dot{\hat z}=A\hat z+B(y+\nu), && \hat x=D(\hat z,\hat c),\ \ \hat z(0)=\hat T_0(\hat x_0).
\end{align*}
For the latent variant, $T_0\in C^1$ is a stationary map with left inverse $T_0^+$ ($\ell_{\rm dec}$-Lipschitz,
$\varepsilon_{\rm dec}=\sup_{\mathcal X}|T_0^+T_0(x)-x|$), the residual along trajectories is
$r=\partial_xT_0(x)f(x,u)-AT_0(x)-Bh(x)-\Phi(T_0(x),c)$ with $|r|\le\varepsilon_{\rm res}$, and $\ell_{\Phi,c}$ is the
Lipschitz constant of $\Phi(z,\cdot)$ \emph{uniformly over} $z\in T_0(\mathcal X)+B(r)$ (the set used in (a)).
$\mu_2(M)=\lambda_{\max}(\tfrac12(M+M^\top))$ is the logarithmic norm.

\begin{proposition}[where the input enters]\label{prop:m1-where}
Under Assumption~\ref{ass:m1}:
\begin{enumerate}
\item[(a)] \emph{Latent injection: stability depends on the input} [proven]. Let
$\bar d:=\varepsilon_{\rm res}+\|B\|\delta_y+\ell_{\Phi,c}\,\delta_c$ and let $r\ge|e(0)|+\bar d/\mu$. If
$\mu_2\big(A+\partial_z\Phi(z,\hat c)\big)\le-\mu<0$ for all $z$ in the set $\mathcal Z_r:=T_0(\mathcal X)+B(r)$ and all
contexts $\hat c$ the observer meets, then $e=\hat z-T_0(x)$ stays in $B(r)$ for all $t$ (so the hypothesis is
checked on a set fixed in advance, not on the solution) and
\[
|e(t)|\le e^{-\mu t}|e(0)|+\tfrac1\mu\big(\varepsilon_{\rm res}+\|B\|\delta_y+\ell_{\Phi,c}\,\delta_c\big),\qquad
|\hat x-x|\le\ell_{\rm dec}|e|+\varepsilon_{\rm dec}.
\]
\item[(b)] \emph{Decoder side: stability is free, extrapolation is not} [proven]. The latent is an input-free LTI filter,
so $\hat z-z$ is driven only by $\nu$ and the initial transient, and $|\hat x-x|$ obeys the bound of
Proposition~\ref{prop:m3-iss} for every input, with the decoder's Jacobian norms $\ell_z(\hat c),\ell_c(\hat c)$ taken at
the contexts actually queried. Both conditioner families have $c$-uniform bounds on a bounded context set, but with very
different constants: for an additive conditioner $\ell_z(c)\le\ell_z(T_0^+)+2\sup\|\partial_zR\|$; for multiplicative
ones (FiLM, LoRA, hypernetwork) the available bound is the product form
$\ell_z(c)\le\prod_l\|W_l+\Delta W_l(c)\|\sup|\sigma'|^{d}$, finite because $\theta(c)$ is bounded (a continuous trunk on a
bounded context set: the SFM context and $g$ have $\tanh$ output layers \vthree{M1-5}{was: tanh trunk; the trunk is SiLU}), but $10^1$--$10^5$ times larger at random initialisation (check [5]) before the tanh base saturates.
\end{enumerate}
\end{proposition}
\noindent\emph{Intuition.} Injecting the input makes the observer a recurrent network whose stability depends on how
strongly the injection reacts to the latent, but its decoder never sees an unfamiliar context; conditioning the decoder
keeps the recurrent part a fixed stable filter, but the network is queried at every unseen context directly.
\begin{proof}
(a) $\dot e=Ae+[\Phi(\hat z,\hat c)-\Phi(T_0x,\hat c)]+d$, $d=B\nu+[\Phi(T_0x,\hat c)-\Phi(T_0x,c)]-r$, $|d|\le\bar d$. While
$e\in B(r)$ the segment $[T_0x,\hat z]$ lies in the ball $T_0x+B(r)\subseteq\mathcal Z_r$ (convexity of the ball,
not of $\mathcal Z_r$, is what is used); write the bracket as $\bar Je$,
$\bar J=\int_0^1\partial_z\Phi(T_0x+se,\hat c)ds$, so $\mu_2(A+\bar J)\le-\mu$ by subadditivity and convexity of $\mu_2$,
and $D^+|e|\le-\mu|e|+\bar d$. The comparison function $\phi(t)=e^{-\mu t}|e(0)|+\bar d(1-e^{-\mu t})/\mu\le r$ bounds
$|e|$ as long as $e\in B(r)$, and $\phi<r$ whenever $|e|=r$ would be reached, so $B(r)$ is forward invariant and the
bound holds for all $t$ (contraction, \citealp{lohmiller1998contraction}; equivalently the perturbation lemmas of
\citealp[Ch.~9]{khalil2002nonlinear}). The state bound is the triangle inequality through $T_0^+(T_0(x))$.
(b) The latent statement is linearity; the bound is Proposition~\ref{prop:m3-iss}; the Jacobian bounds are the chain rule
($R$'s $z$-Jacobian does not involve $c$ multiplicatively, whereas modulated layers multiply $z$-features by $c$-dependent
weights), and $\theta(c)$ is bounded because the trunk is continuous and its input is bounded ($\tanh$ output layers of
$\psi$ and $g$). For the default, $\hat z$ is replaced by the stack $Z$ in the Jacobians (Proposition~\ref{prop:m3-iss}).
\end{proof}

\begin{corollary}[the certified margin shrinks with the context]\label{cor:m1-amplitude}
$\mu_2(A+\partial_z\Phi)\le\mu_2(A)+\ell_\Phi$, and $\mu_2(A)=-\alpha$ for normal $A$, so $\mu\ge\alpha-\ell_\Phi$
(for general $A$, the Gronwall version gives $\alpha-\kappa\ell_\Phi$). For an injection with exact zero-input
collapse, $\Phi(z,c)=M(z,c)-M(z,0)$, $M\in C^2$: $\ell_\Phi(c)\le\min\big(|c|\,L_{zc},\,2\sup\|\partial_zM\|\big)$ with
$L_{zc}=\sup\|\partial_c\partial_zM\|$, and $|c|\le\ell_E\|u\|_\infty$ (sup norm) since $E(0)=0$; for an input-affine
injection $\Omega(z)u$, $\ell_\Phi\le\ell_\Omega\|u\|_\infty$. The margin is evaluated at the \emph{observer's} context
$\hat c$, so a corrupted input shrinks it. The margin saturates: if $2\kappa\sup\|\partial_zM\|<\alpha$ the latent observer
is certified at every amplitude, and when the plant input saturates (as through a saturating actuator) the exact
correction's $z$-Lipschitz constant stays bounded too. The condition is \emph{sufficient only}: for the code's trained
injection the certified margin is already negative at moderate gain ($-5.47$, log-norm $+8.1$ in check [3]) while two
observer copies still contract, so a negative margin predicts nothing, and only a measured incremental rate does.
\end{corollary}

\begin{proposition}[the margin cannot be removed at a fixed spectrum]\label{prop:m1-tight}
[proven] There are bounded, observable plants with an \emph{exact} injection ($\varepsilon_{\rm res}=0$) whose latent
observer diverges while a decoder-side observer with a bounded sufficient context is within $\tfrac{\ell}{2}\varepsilon(L)$
of exact. For $\dot x=(u-a)x+b$, $y=x$, $a>\alpha>0$, scalar latent $\dot z=-\alpha z+y$ and constant $u\in(\alpha,a)$:
the plant converges to $b/(a-u)$, the stationary map is $T_0(x)=Px+q$, $P=1/(\alpha-a)$, $q=-Pb/\alpha$, the exact
injection is $\Phi^*(z,u)=u(z-q)$ ($\ell_\Phi=|u|$), and $\dot e=(u-\alpha)e$. The windowed map is
$T_L(x,w)=P_L(u)\,x+Q_L(u)$ with $P_L=\int_{-L}^0e^{(\alpha+u-a)\tau}d\tau>0$ and
$Q_L=\int_{-L}^0e^{\alpha\tau}\,b\,\tfrac{e^{(u-a)\tau}-1}{u-a}\,d\tau$ (both bounded functions of the window input), so the
decoder $\hat x=(\hat z-Q_L)/P_L$ errs by at most $\tfrac{\ell}{2}\varepsilon(L)$ with $\ell=1/P_L$
(Lemma~\ref{lem:m1-window}), for every $u\in(\alpha,a)$.
\end{proposition}
\begin{proof}
$T_0'(x)f(x,0)=P(b-ax)$ and $-\alpha T_0+x=(1-\alpha P)x-\alpha q=-aPx+Pb$; so $T_0$ solves the zero-input PDE and
$\partial_xT_0[f(x,u)-f(x,0)]=Pux=u(T_0(x)-q)$. Then $\dot e=-\alpha\hat z+x+u(\hat z-q)-P((u-a)x+b)=(u-\alpha)e$.
For the window: $X(x,0;\tau;u)=e^{(u-a)\tau}x+b(e^{(u-a)\tau}-1)/(u-a)$ in \eqref{eq:m1-TL}.
\end{proof}
\noindent Here the small-gain threshold $u=\alpha$ is attained exactly. The \emph{exact} time-varying map $z=m(t)x+n(t)$,
$\dot m=1-(\alpha+u-a)m$, $\dot n=-\alpha n-bm$, is also a sufficient context, but for $u<a-\alpha$ its coefficient
$-(\alpha+u-a)$ is positive and $m,n\to\infty$ (check: $m(30)=2\cdot10^{19}$ at $u=0.5$); it is outside the bounded
class of Proposition~\ref{prop:m1-where}(b), which is why the windowed context is the one used. The obstruction is
intrinsic to the stationary route \emph{at a fixed spectrum}: a designer removes it by choosing $\alpha>\bar u$ (here) or
$\alpha>\bar u\ell_\Omega\kappa$ (Corollary~\ref{cor:m1-amplitude}), i.e.\ by a gain that grows with the input bound, which
is also what \citet[Theorem~4]{bernard2019luenberger} require. A second example, the rotating oscillator
$\dot x=(\omega_0+u)Jx$, $y=x_1$ ($|x|$ is conserved for every input), with the drift map $P$ and the exact correction
$\Omega(z)=PJP^+z$, has error rate $\max\operatorname{Re}\lambda(A+uPJP^+)$: stable for all tested $u>-\omega_0$, but
unstable ($+2.65$, $+5.16$ at $u=-2,-4$, $\omega_0=1$, $n_z=5$, $\alpha=1.5$) when the rotation reverses, although the
plant is bounded and observable there; this is the input-independence required by the stationary route
\citep[Theorem~4]{bernard2019luenberger}, which also needs an input-affine plant and a gain growing with $\bar u$.
The decoder-side route needs only Assumption~\ref{ass:m1} and a sufficient context.

\paragraph{What the bounds predict, and what the current nodes test.} Under protocol v2 no tier changes the input
amplitude (every family lives in the same box), so the quantity that varies between tiers is the \emph{distance in context
space}, $d(\hat c,\mathcal C_{\rm train})$, not $|c|$. (P1) The latent variant's excess error on \texttt{shape\_ood} and
\texttt{sequence\_ood} is the extrapolation error of $\Phi$ at the new contexts, filtered at the measured incremental rate;
it degrades sharply only if that rate approaches zero. Decoder-side variants pay $\ell_z(\hat c)\,|\Delta z|$ and
$\ell_c(\hat c)\,|\Delta c|$ at the new contexts (Corollary~\ref{cor:m1-ood}), with the additive conditioner's $\ell_z$ small
and $c$-uniform. (P2) Input noise moves the latent variant's operating point to a corrupted context $\hat c$ (its margin is
evaluated there), whereas decoder-side variants stay stable for every input and pay $\ell_c\ell_E\delta_u$.
\emph{Status: untested.} No current node measures $\mu_2$ or the incremental rate on visited segments, $\ell_z(\hat c)$,
$\ell_c(\hat c)$, or the context distance per tier; under C3-faults \texttt{input\_noise} every variant, latent included,
degrades by about the same factor ($\approx10\times$ on vdp, read from the C3-faults node report; not a number of the
generated results layer), which neither supports nor falsifies (P2). The quantities to
record are eval-only on saved checkpoints: $\min_{c'\in\mathcal C_{\rm train}}\|\hat c-c'\|$ per tier (C1-where,
TEST-C1), the latent margin $\alpha-\kappa\ell_\Phi$ already computed by \texttt{metrics.py}, and the realised gains of
Definition~\ref{def:m3-measured}. Measured Lipschitz constants and log-norms are estimates on the visited set (lower
bounds of suprema), not certificates.

\begin{corollary}[out of distribution: Lipschitz extrapolation, and nothing more]\label{cor:m1-ood}
[proven] Let $\varepsilon_{\rm train}:=\sup|D(z',E(w'))-x'|$ over the training support $S_{\rm train}$ of triples
(label latent $z'=T^0$, window $w'$, state $x'$), and \emph{assume} the target $x=G^*(z,w)$ of Lemma~\ref{lem:m1-window}
has, on its domain $S_L$, Lipschitz constants $\ell$ in $z$ (this is (A4)) and $\ell_w$ in $w$ (sup norm; a new assumption,
not implied by sufficiency), and assume its extension beyond $S_L$ keeps these constants (a componentwise McShane
extension achieves them only up to a factor $\sqrt{n_x}$, which then multiplies $\ell$ and $\ell_w$ in (a)).
(a) For a query $(z,w)$ and any $(z',w')\in S_{\rm train}$, with $D$'s Jacobian norms $\ell^D_z,\ell^D_c$ on the segment,
\[
|D(z,E(w))-G^*(z,w)|\le\varepsilon_{\rm train}+\tfrac{\ell}{2}\varepsilon(L)+(\ell^D_z+\ell)|z-z'|+(\ell^D_c\,\ell_E+\ell_w)\|w-w'\|_\infty ,
\]
the $\tfrac{\ell}{2}\varepsilon(L)$ because $G^*(z',w')=x'$ only at $z'=T_L$, whereas the training latent is $T^0$; and
$|G^*(z(t),u_t)-x(t)|\le\tfrac{\ell}{2}\varepsilon(L)$ (Lemma~\ref{lem:m1-window}).
\vthree{M1-6}{stack} For the default, replace $z$ by the stack $Z$, $|z-z'|$ by $\|Z-Z'\|$ and $\ell$ by the Lipschitz constant of
the target in $Z$; the proof is the same triangle inequality.
(b) The linear rate cannot be improved without structure. Among targets that are $2\lambda$-Lipschitz and agree with a
$\lambda$-Lipschitz one on $S_{\rm train}$, $G_\pm=G_0\pm\lambda\,d(\cdot,S_{\rm train})\,e$ ($|e|=1$) both fit the data,
so every learner errs by at least $\lambda\,d(q,S_{\rm train})$ at some query $q$. At the plant level, if $\varphi$ is a
diffeomorphism of $\mathcal X$ with $h\circ\varphi=h$, equal to the identity on an open neighbourhood of the states reached
under the training inputs (so that $D\varphi=I$ there), then the plants $f$ and $\varphi_*f$ produce identical training
data, and under an input that drives the state where $\varphi\neq\mathrm{id}$ they produce identical $(y,u)$ from $x_0$ and
$\varphi(x_0)$, so any observer errs by at least $\tfrac12|\varphi(x(t))-x(t)|$ on one of them.
\end{corollary}
\begin{proof}
(a) Triangle inequality through $D(z',E(w'))$, $G^*(z',w')$ and $G^*(T_L(x',w'),w')=x'$, with
$|E(w)-E(w')|\le\ell_E\|w-w'\|_\infty$ and $|z'-T_L(x',w')|\le\tfrac12\varepsilon(L)$.
(b) $d(\cdot,S)$ is 1-Lipschitz and vanishes on $S$; $|G_+-G_-|=2\lambda d$ and one learner output is at distance
$\ge\lambda d$ from one of them. For the plants: $\varphi$ maps solutions of $f$ to solutions of $\varphi_*f$ and preserves
$y$; on the training set it is the identity.
\end{proof}
\noindent\emph{What it predicts under protocol v2.} Distances are measured in input (window) space so that they are
comparable across encoders. In v2 at most $0.1\%$ of evaluation states lie outside the training box on any dev system
and tier (C1-where node report; not in the generated results layer), and the
input amplitude box is shared by all tiers, so $|z-z'|$ and the state-level construction of (b) are inert: the relevant
distance is $\|w-w'\|_\infty$ (equivalently $d(\hat c,\mathcal C_{\rm train})$ in context space), which \texttt{shape\_ood}
and \texttt{sequence\_ood} change through the waveform and the event timing at the same amplitude. The corollary therefore
predicts a shape/sequence-OOD excess that is at most linear in that context distance and, by (b), can be no better
without structure; it does not predict its size, and the context distance per tier is not yet reported (eval-only; C1-where,
TEST-C1). A model-based filter with the correct model extrapolates by structure; a learned observer cannot, beyond this
Lipschitz rate.

\subsection{The autonomous case}

\begin{proposition}[no input: a fixed filter and a stationary decoder of the latent window]\label{prop:m1-autonomous}
\vthree{M1-7}{new; replaces the zero-input paragraph, which claimed collapse to the autonomous KKL observer and holds
only for the C1 arms} [(a), (b), (c) proven; (d) is what transfers, with its hypotheses] Let $u\equiv0$ (an autonomous
plant, or an input held at zero) and $\xi_0=0$.
\begin{enumerate}[label=(\alph*)]
\item (\emph{exact reduction}) $c_k=0$ for every $k$, and the deployed observer is
$\hat x_k=D_0(Z_k)$ with $D_0(Z):=T^+(z_0;[0,g(Z)])$ ($z_0$ the first block of $Z$): the fixed input-free latent filter
followed by a time-invariant map of the latent window. $D_0$ is in general neither the phase-1 decoder $T_0^+$ nor the
retrained base: $\Delta W_l(\mathrm{Trunk}([0,g(Z)]))$ vanishes where $g(Z)=0$ and not in general elsewhere.
\item (\emph{latent window along the solution}) Suppose there is a continuous $T:\mathcal X\to\mathbb R^{n_z}$ with
$\tfrac{d}{dt}T(x(t))=AT(x(t))+Bh(x(t))$ along every solution in $\mathcal X$ (the autonomous KKL map; conditions for its
existence and injectivity: \citealp[Theorems~2.2, 2.4]{andrieu2006existence}, \citealp[Theorem~3.4]{brivadis2023further}).
For $t_k\ge S$ let $\mathcal T(x(t_k)):=(T(x(t_k)),T(x(t_k-s\,dt)),\dots,T(x(t_k-S)))$; the entries are values of
$T$ along the true past, which for an autonomous $C^1$ flow is a function of $x(t_k)$. Then, with
$e_0:=\hat z_0-T(x_0)$, noise-free measurements and $r_K(t):=(\sum_{j<K}e^{-2\alpha(t-js\,dt)})^{1/2}$,
\[
\|Z_k-\mathcal T(x(t_k))\|\le\kappa\,r_K(t_k)\,\|e_0\|+\sqrt K\,e_{\rm disc}.
\]
\item (\emph{error bound}) Let $\varepsilon_0:=\sup|D_0(\mathcal T(x))-x|$ over the states $x(t)$, $t\ge S$, of the
admissible solutions, and let $D_0$ be $\ell_Z$-Lipschitz on a set containing the segments
$[\mathcal T(x(t_k)),Z_k]$. Then for $t_k\ge S$
\[
|\hat x_k-x(t_k)|\le\varepsilon_0+\ell_Z\big(\kappa\,r_K(t_k)\,\|e_0\|+\sqrt K\,e_{\rm disc}\big),
\]
and $r_K(t)\le e^{-\alpha(t-S)}/\sqrt{1-e^{-2\alpha s\,dt}}$. Sensor deviations add the terms of
Proposition~\ref{prop:m3-iss} with $\Delta c=0$.
\item (\emph{what transfers from autonomous KKL theory}) Where $T$ exists and is injective with a continuous left
inverse $T^*$ (the cited conditions), $\mathcal T$ is injective (its first block is) and $T^*\circ\pi_1$ is a stationary
left inverse of $\mathcal T$ in the class of $D_0$; so the autonomous existence results apply to the target of $D_0$,
and (c) has the form of the robustness bound of learned KKL observers \citep[Lemma~17 and Theorem~19 of arXiv v2;
Lemma~6 and Theorem~4 of v3]{niazi2025kkl} (their $\mathrm{cond}(V)$ is our $\kappa$), with the stack in place of the latent. What does
not transfer is any a priori value of $\varepsilon_0$ or $\ell_Z$: the trained $D_0$ is not $T^*\circ\pi_1$, and
$\varepsilon_0$ is a measured quantity, exactly as in the input-driven case.
\end{enumerate}
\end{proposition}
\emph{Intuition.} With no input the SFM context is identically zero, so nothing in the observer depends on time
except the latent path; what remains is a KKL observer whose decoder reads a window of its own latent.
\begin{proof}
(a) $\xi_k=\sum_{l\ge1}(A^\xi_d)^{l-1}B_du_{k-l}/u_{\rm ref}=0$; $\psi$ is bias-free with $\tanh(0)=0$ in every layer, so
$c_k=\psi(0)=0$. The weights are fixed at deployment, so $D_0$ does not depend on $k$. The trunk receives $[0,g(Z)]$;
$V_l$ is linear and bias-free, so $\Delta W_l=0$ when $\mathrm{Trunk}([0,g(Z)])=0$, which holds at $g(Z)=0$ (bias-free
trunk) and not in general otherwise. (b) With noise-free $y=h(x)$, $e(t):=z(t)-T(x(t))$ obeys $\dot e=Ae$ for the
continuous latent, so $\|e(t)\|\le\kappa e^{-\alpha t}\|e_0\|$; for $t_k\ge S$ every stack entry is at a time
$t_k-js\,dt\ge0$, and the Euclidean norm of the stack is the root sum of squares. The deployed latent is the exact
step of Lemma~\ref{lem:D}(a) with a piecewise-linear $y$; its difference to the continuous latent from the same initial
condition does not depend on that condition (linearity) and is at most $e_{\rm disc}$ per entry (Lemma~\ref{lem:D}(b)).
(c) $|\hat x_k-x(t_k)|\le|D_0(Z_k)-D_0(\mathcal T(x(t_k)))|+|D_0(\mathcal T(x(t_k)))-x(t_k)|$, then (b); the bound on
$r_K$ is the geometric series. (d) $T^*(\pi_1\mathcal T(x))=T^*(T(x))=x$.
\end{proof}
\noindent\emph{Status of the C1 zero-input claim.} The statement ``with zero input HyperKKL collapses exactly to the
autonomous KKL observer'' holds for the C1 arms ($K=1$, frozen base, $E(0)=0$; \texttt{tests/test\_core.py}) and is
\emph{false} for the default: on the stored default models (battery and cartpole, two seeds each) the context is exactly
$0$ at $u=0$ while $g(Z)\ne0$ and $\|\Delta W_l\|_2$ is $0.2$ to $204$ on random stack probes (check V7b; off-data
probes, so these magnitudes are not on-data values). \emph{Check} V6: linear autonomous plant, perturbed stationary
decoder of the stack, bound (c) holds (sup ratio $0.04$).

\subsection{The non-autonomous residual (optional; not in the default)}

\begin{remark}[the correct physics residual for an input-driven decoder]\label{rem:m1-tv}
\vthree{M1-8}{new; replaces nothing, but no statement in v3 says that a physics term helps} For the autonomous phase-1
map the KKL equation is $\partial_xT_0\,f(x,0)=AT_0+Bh$. For an input-driven observer the object that is deployed is the
decoder, and the equation it must satisfy is the \emph{inverse} non-autonomous equation: if $D(Z(t),c(t))=x(t)$ along
trajectories (an exact left inverse), then between samples (input held at $u_k$)
\[
R:=\sum_{j=0}^{K-1}\partial_{z^{(j)}}D\,\big(Az^{(j)}+By(t-js\,dt)\big)+\partial_cD\,\dot c-f\big(D(Z,c),u_k\big)=0,\qquad
\dot c=\partial_\xi\psi\,(A^\xi\xi+B^\xi u_k/u_{\rm ref}),
\]
where $z^{(j)}$ is the $j$-th stack entry, each obeying the latent ODE, and the direct path $u_k/u_{\rm ref}$ has zero
derivative on the step. [proven: differentiate the identity with the chain rule; $f$ is the true right-hand side.] Both
chain-rule terms are needed: for an exact linear decoder the residual is $0$ to rounding with them and $O(10^{-1})$ if
either the $\dot c$ term or the delayed-latent term is dropped (check V5, $K=1$ and a two-entry FIR decoder). This is the
form implemented in \texttt{experiments/tv\_physics} (\texttt{tv.py}). The term is \emph{optional} and is not part of
the default. It needs a model of $f$; with a wrong model its zero is the model's, not the plant's. Evidence (recipe
A4+D3, not the default, three seeds): no significant effect in $10/10$ cells, also with the true $f$
(\texttt{tv\_physics/RESULTS.md}); the phase-1 autonomous PDE term was neutral or harmful there. We make no claim that
either physics term improves the learned observer.
\end{remark}

\paragraph{Numerical checks.} \vthree{M1-9}{v3 checks added} \texttt{theory/checks/v3\_checks.py}: V5 (residual of
Remark~\ref{rem:m1-tv}), V6 and V7 (Proposition~\ref{prop:m1-autonomous}), V8 (stride rule). \texttt{theory/checks/m1\_location.py}
(written for the $K=1$ arms): [1] zero-input collapse of all encoders and
conditioners (errors $\le5\cdot10^{-16}$); [2] at fixed $c$, FiLM, LoRA and hypernetwork decoders equal base-architecture
networks with explicit weights ($\le3\cdot10^{-13}$); [3] the injection's $z$-Lipschitz constant is $0,0.61,1.22,2.33,4.26,6.82,7.29$ at
$|c|=0,0.25,0.5,1,2,8,32$ (linear, then saturating), and for two observer copies the measured incremental rate is below the
sup log-norm on the visited segments ($-1.52\le-1.23$; $-1.05\le-0.29$, where the small-gain margin is $-5.47$ and certifies
nothing; at the largest gain the log-norm is $+8.1$ and also certifies nothing while the copies still contract, so the
log-norm condition is sufficient only); [4] the rotating oscillator; [5] \emph{at random initialisation} (trained
decoders are not tested here) the $z$-Jacobian of concat stays below its uniform-in-$c$ bound ($\le1.25$ against $12$ for
$|c|\le256$) while FiLM, LoRA and hypernetwork reach $81$, $3.7\cdot10^5$, $2.3\cdot10^6$ and then fall (LoRA and
hyper to $0$ and $2.9$ at $|c|=256$): $c$-uniformly bounded, with a far larger constant, as stated in
Proposition~\ref{prop:m1-where}(b); [6] the linear windowed map is affine with $w$-independent slope to $10^{-15}$;
[7] the Koopman resolvent identity $p_i=b_iC(F-\lambda_iI)^{-1}$ (to $10^{-16}$) and its Laplace-transform form (to
$10^{-8}$); [8] the a-priori form of Proposition~\ref{prop:m1-where}(a): with $\mu=0.46$, $|e(0)|=0.5$, $\sup|d|=0.3$ the
error never leaves $B(r)$, $r=1.15$. \texttt{theory/checks/check\_theory.py}: \texttt{where\_input\_enters}
(Proposition~\ref{prop:m1-tight}: measured rate $u-\alpha$ to $10^{-3}$; windowed-context error equals
$\tfrac{\ell}{2}\varepsilon(L)$ to three digits at $L=6$, i.e.\ the bound is attained for a constant input; the exact
context's $m(t)$ reaches $10^{11}$--$10^{19}$).

% ---- bibliography entries used (for the main .bib) ----
% bernard2019luenberger: P. Bernard, V. Andrieu. Luenberger observers for nonautonomous nonlinear systems.
%   IEEE TAC 64(1):270--281, 2019. doi:10.1109/TAC.2018.2872202  (numbering checked on HAL hal-02050020)
% andrieu2006existence: V. Andrieu, L. Praly. On the existence of a Kazantzis--Kravaris/Luenberger observer.
%   SIAM J. Control Optim. 45(2):432--456, 2006.  (numbering checked on HAL hal-00362726)
% brivadis2023further: L. Brivadis, V. Andrieu, P. Bernard, U. Serres. Further remarks on KKL observers.
%   Systems & Control Letters 172:105429, 2023.  (numbering checked on HAL hal-03695863)
% niazi2025kkl: M.U.B. Niazi, J. Cao, M. Barreau, K.H. Johansson. KKL observer synthesis for nonlinear systems via
%   physics-informed learning. arXiv:2501.11655v2, 2025.
% boyd1985fading: S. Boyd, L.O. Chua. Fading memory and the problem of approximating nonlinear operators with
%   Volterra series. IEEE Trans. Circuits Syst. 32(11):1150--1161, 1985.
% leshno1993multilayer: M. Leshno, V.Ya. Lin, A. Pinkus, S. Schocken. Multilayer feedforward networks with a
%   nonpolynomial activation function can approximate any function. Neural Networks 6(6):861--867, 1993.
% zeng2024expressive: Y. Zeng, K. Lee. The expressive power of low-rank adaptation. ICLR 2024.
% lohmiller1998contraction: W. Lohmiller, J.-J.E. Slotine. On contraction analysis for non-linear systems.
%   Automatica 34(6):683--696, 1998.
% khalil2002nonlinear: H.K. Khalil. Nonlinear Systems, 3rd ed. Prentice Hall, 2002.
% korda2018linear: M. Korda, I. Mezic. Linear predictors for nonlinear dynamical systems: Koopman operator meets
%   model predictive control. Automatica 93:149--160, 2018.
% proctor2016dmdc: J.L. Proctor, S.L. Brunton, J.N. Kutz. Dynamic mode decomposition with control.
%   SIAM J. Applied Dynamical Systems 15(1):142--161, 2016.
